\documentclass[12pt]{article}
\usepackage[utf8]{inputenc}
\usepackage[T1]{fontenc}
\usepackage{hyperref}
\title{Ngunawal Language Resources for K-12 Curriculum Development}
\author{Compiled by Claude Sonnet 5.5 (claude-sonnet-5-5) for Andr\'as Kornai}
\date{}
\begin{document}
\maketitle
\section{Introduction}
Ngunawal (also spelled Ngunnawal; ISO 639-3 xul) is a sleeping language of the Canberra region of Australia. Documentation is meagre: mostly early manuscript sources and only a few minutes of audio. Revival work is led by community groups with AIATSIS. Little is publicly available online, and no published dictionary was found in the search results.
\section{Primary Educational Resources}
\begin{itemize}
\item AIATSIS news item on a pilot language program at Fraser Primary School, developed by community members and AIATSIS linguists, with a standardised orthography: \url{https://aiatsis.gov.au/whats-new/news/school-children-benefit-language-revival}
\end{itemize}
\section{Linguistic and Reference Materials}
\begin{itemize}
\item AIATSIS Ngunnawal Language Revival Project. A linguist analysed historical materials and wrote a technical grammar, which the community uses to build learning materials: \url{https://aiatsis.gov.au/research/current-projects/ngunnawal-language-revival-project}
\item AIATSIS publication record: \url{https://aiatsis.gov.au/publication/116700}
\item ``First steps on the Ngunawal language revitalization journey'' (academic paper): \url{https://core.ac.uk/works/19007414}
\item Scholarly repository record found in the same search: \url{https://scholarspace.manoa.hawaii.edu/handle/10125/25315}
\end{itemize}
\section{Religious and Community Resources}
Community groups named in the sources are the Ngaiyuriija Ngunawal Language Group and, since 2018, the Winanggaay Ngunnawal Language Group. No contact details were found in the search results. The University of Canberra event page mentions a linguist working with community members on a dictionary: \url{https://canberra.edu.au/events/Home/Event/461692}
\section{Digital and Audio Resources}
No public audio or digital database was found in the search results. Historical audio is described as only a few minutes.
\section{Curriculum Development Context}
A school pilot exists, so a curriculum can build on it, working in partnership with the community language groups and AIATSIS. Because the source material is thin, content would mostly be vocabulary, greetings, place names and cultural learning, developed under community authority. The status of the dictionary and grammar materials should be confirmed with AIATSIS and the community groups.
\end{document}